\documentclass[
  a4paper,
  12pt,
  oneside,
  answers=true,
  solutions=inline
]{maki-notes}

\renewcommand{\LectureNotesTitle}{Inline Solutions Test}
\renewcommand{\AuthorInfo}{Test Runner}

\begin{document}
\mainmatter

\chapter{Inline Answers}

\section{Exercise Flow}

\begin{exercise}[等差数列求和]
\label{ex:inline-sum}
设 \(a_n = 2n-1\). 求 \(\sum_{n=1}^{5} a_n\).
\end{exercise}

\begin{solution}
这是前五个正奇数之和, 因此
\[
  \sum_{n=1}^{5}(2n-1)=5^2=25.
\]
\end{solution}

上面的答案应当以内联方式显示, 且可以通过 \exerciseref{ex:inline-sum} 引用.

\begin{center}
\MakiFitTikZ[0.35\linewidth]{%
\begin{tikzpicture}[x=1cm,y=1cm]
  \node[flow terminal] (a) at (0,0) {题目};
  \node[flow process] (b) at (2.6,0) {思路};
  \node[flow terminal] (c) at (5.2,0) {答案};
  \draw[flow arrow] (a) -- (b);
  \draw[flow arrow] (b) -- (c);
\end{tikzpicture}%
}
\end{center}

\end{document}
